\documentclass[11pt]{article}
\usepackage[utf8]{inputenc}
\usepackage[T1]{fontenc}
\usepackage[margin=1in]{geometry}
\usepackage{hyperref}

\title{Previously Published / Preprint Statement}
\author{Hemanth Manchabale Papachappa \and Aniriuddha Ganguly}
\date{September 2026}

\begin{document}
\maketitle

\noindent\textbf{Manuscript title:} Geodesic Diagnostic Trajectories (GDT):
A Differential Geometry Framework for Adaptive Decision Support and Operational
Optimization in EHR Systems

\bigskip

\noindent The authors confirm that this manuscript has \textbf{not} been published
in whole or in part in any peer-reviewed journal and is \textbf{not} currently
under review at any other journal.

\bigskip
\section*{Related Prior Work by the Authors}

The following prior publications and preprints are related to but substantively
distinct from the present submission:

\begin{enumerate}

\item \textbf{Manchabale Papachappa, H., Kumar, A., and Badam, N.}
``Riemannian Intent Manifolds for Session-Based Search: A Joint Embedding
Predictive Architecture for Intent Forecasting.'' 2026 (preprint).\\
\textit{Relationship:} This preprint introduces the general JEPA-based intent
forecasting framework on Riemannian manifolds for web-scale session search. The
present GDT manuscript extends that work to the clinical EHR domain, adds the
geodesic shift interference gate, proves Proposition~1 (context interference
bound), and reports a four-arm clinical benchmark on MIMIC-III/IV, eICU, and
HiRID. The clinical benchmark, formal interference bound, confidence-scaling
mechanism, and queueing model are new contributions not present in the preprint.

\item \textbf{Manchabale Papachappa, H.}
``Continuum Memory Architecture (CMA) for Clinical Search: A Simulation-Based
Evaluation of Intent-Aware Retrieval Using Complex Public Clinical Data.''
SSRN preprint, DOI: 10.2139/ssrn.7197862, 2026.\\
\textit{Relationship:} CMA is a direct predecessor of GDT and serves as the
primary comparison system in the four-arm benchmark reported here. GDT differs
from CMA in three concrete ways: (1) GDT adds a confidence-scaling mechanism to
the JEPA prefetch engine (CMA uses fixed-weight prefetch); (2) GDT formalizes and
proves the context interference bound as Proposition~1 (CMA does not); (3) GDT
introduces the M/M/c queueing model for operational throughput projection (not
present in CMA). These differences are fully described and quantified in the
manuscript.

\end{enumerate}

\bigskip
\section*{Preprint Policy Compliance}

Per JTEHM policy, authors are permitted to post preprints prior to submission.
Should the authors post a preprint version of this manuscript to arXiv or a
comparable server during the review period, the journal will be notified immediately.
The preprint will clearly identify itself as a manuscript under review and will not
be promoted in a manner that could compromise the integrity of the peer-review process.

\bigskip
\noindent\rule{\textwidth}{0.4pt}

\bigskip
\noindent\textbf{Author signatures:}

\bigskip
\noindent\begin{tabular}{ll}
Hemanth Manchabale Papachappa & \underline{\hspace{6cm}} \\[12pt]
Aniriuddha Ganguly             & \underline{\hspace{6cm}} \\
\end{tabular}

\bigskip
\noindent\textbf{Date:} September 2026

\end{document}
